\chapter{Introdução}
\label{cap:introducao}

\orientacao{Este capítulo apresenta o projeto de forma geral, estabelecendo seus objetivos, sua justificativa, seu escopo e a organização do documento. Sua função é fornecer ao leitor uma visão inicial que permita compreender o propósito e a abrangência do trabalho antes de entrar no contexto e nos aspectos técnicos.}

\section{Caracterização do Projeto}
\orientacao{Apresentar uma descrição geral do projeto, caracterizando sua natureza, finalidade, origem e principais características, de modo a situar o leitor antes da apresentação dos objetivos, da justificativa e do escopo.}

\section{Objetivos}
\orientacao{Apresentar os objetivos que orientam o desenvolvimento do projeto e delimitam os resultados que se pretende alcançar.}

\subsection{Objetivo Geral}
\orientacao{Descrever, de forma ampla, o principal resultado pretendido com o desenvolvimento do projeto.}

\subsection{Objetivos Específicos}
\orientacao{Relacionar os resultados parciais ou etapas necessárias para alcançar o objetivo geral.}

\section{Justificativa}
\orientacao{Apresentar as razões que justificam a realização do projeto, considerando sua relevância para o público atendido, para a comunidade e para a formação dos estudantes.}

\section{Escopo do Projeto}
\orientacao{Delimitar o que faz parte do projeto e, quando necessário, o que está explicitamente fora de seu escopo.}

\section{Estrutura do Documento}
\orientacao{Apresentar brevemente a organização do documento do projeto, indicando o conteúdo e a finalidade de suas principais seções.}
